% Generated draft: requires mathematical and source review.
\documentclass[11pt]{article}
\usepackage[margin=25mm]{geometry}
\usepackage{fontspec}
\setmainfont{Latin Modern Roman}
\usepackage{amsmath,amssymb,mathtools,unicode-math}
\setmathfont{Latin Modern Math}
\usepackage{enumitem,needspace,xurl,hyperref}
\setlength{\parindent}{0pt}
\setlength{\parskip}{5pt}
\setlength{\emergencystretch}{3em}
\newcommand{\N}{\mathbb{N}}
\newcommand{\Z}{\mathbb{Z}}
\newcommand{\Q}{\mathbb{Q}}
\newcommand{\R}{\mathbb{R}}
\newcommand{\C}{\mathbb{C}}
\newcommand{\F}{\mathbb{F}}
\newcommand{\A}{\mathbb{A}}
\newcommand{\PP}{\mathbb{P}}
\newcommand{\E}{\mathbb{E}}
\newcommand{\eps}{\varepsilon}
\newcommand{\dd}{\,\mathrm{d}}
\begin{document}
\section*{UnsolvedMath ID: 20003314}
\textbf{Source problem code:} catalogue entry
\subsection{Definitions and mathematical statement}
The record is only ``Conjecture 32 is related to the following conjecture.'' The words ``OPEN PROBLEMS 14'' are a running page header. Neither neighboring conjecture is stated here, so this is context rather than a new conjecture. The aid supplies a weighted-approximation/diagonal-flow interpretation and a reported status result for neighboring Conjecture 33, without reproducing the full adjacent statements.

Let $\|x\|$ be nearest-integer distance, and define for $0\leq s\leq1$
\[
 C_s=\{(\alpha,\beta):
 \inf_{q\geq1}\max\{q^s\|q\alpha\|,
                         q^{1-s}\|q\beta\|\}>0\},
 \qquad
 \operatorname{Bad}=\{\theta:\inf_{q\geq1}q\|q\theta\|>0\}.
\]
The aid explicitly includes the endpoint identities
\[
 C_0=\R\times\operatorname{Bad},\qquad
 C_1=\operatorname{Bad}\times\R.
\]
They are not obtained by silently discarding a zero-weight coordinate in the displayed definition; the source reports an endpoint-safe proof.

In $\mathcal X_3=SL(3,\R)/SL(3,\Z)$, write
\[
 u_{\alpha,\beta}=\begin{pmatrix}1&0&\alpha\\0&1&\beta\\0&0&1\end{pmatrix},
 \qquad L_{\alpha,\beta}=u_{\alpha,\beta}\Z^3.
\]
For $a,b\geq0$ with $a+b>0$, let
\[
 R_{a,b}=\{\operatorname{diag}(e^{at},e^{bt},e^{-(a+b)t}):t\geq0\}.
\]
Bounded means relatively compact orbit in $\mathcal X_3$. The reported weighted Dani correspondence is
\[
 R_{a,b}L_{\alpha,\beta}\text{ bounded}
 \quad\Longleftrightarrow\quad
 (\alpha,\beta)\in C_{a/(a+b)}.
\]
Boundary rays with one coefficient zero are included; the zero direction is not a ray and would make the ratio undefined.

Consequently, for two such rays, any point of
\[
 C_{a/(a+b)}\cap C_{c/(c+d)},\qquad c,d\geq0,\ c+d>0,
\]
gives a lattice with both prescribed forward ray orbits bounded. Its orbit under the \emph{full} positive diagonal group $A$ is unbounded: the fixed lattice vector $e_1$ is shortened to zero by $\operatorname{diag}(e^{-t},e^t,1)\in A$. This is a conditional construction from an intersection point, not by itself a proof that every required intersection is nonempty. The distinction between a closed forward cone and the full two-sided diagonal group is essential.

The aid separately reports a dimension argument for the neighboring conjecture: the An--Guan--Kleinbock common bounded-orbit set for countably many applicable diagonal rays has Hausdorff dimension eight in the eight-dimensional space $\mathcal X_3$, whereas the full-$A$-bounded set of Einsiedler--Katok--Lindenstrauss has dimension at most two. Hence the former is not contained in the latter, giving lattices bounded along all those rays but not along $A$. The exact ray class and theorem hypotheses are not supplied, so this is retained as an attributed status result, not extended to arbitrary flows or claimed as a solution of an unstated Conjecture 32. The canonical record itself still contains no standalone mathematical target.
\subsection{Short English statement}
This is a transition sentence. The aid connects weighted badly approximable pairs, including endpoint weights, to bounded forward diagonal rays and distinguishes them from a full diagonal-group orbit. Its reported countable-ray existence argument needs the exact external theorem hypotheses and does not replace the missing neighboring conjectures.
\subsection{Sources}
\begin{itemize}
\item Local UnsolvedMath registry, permanent ID 20003314, source code catalogue entry; \texttt{aim-other-notes.json}, notes[120], tagged conjecture despite being a transition.
\item AIM, \emph{Emerging applications of measure rigidity}, bridge after Conjecture 32: \url{https://aimath.org/WWN/measrigid/measrigid.pdf}. ``OPEN PROBLEMS 14'' is a running header, not problem text.
\item Weighted Dani correspondence, An--Guan--Kleinbock common bounded-orbit results and Einsiedler--Katok--Lindenstrauss full-diagonal dimension bounds are named in the supplied aid without full bibliographic entries or hypotheses.
\item Machine-generated research attempt 1: context only, not a problem at the run, rigor labelled ``proved,'' low novelty confidence; registry open/open-triage. Not peer reviewed; independently verify before citation. Registry updated 14 August 2026.
\end{itemize}
\end{document}
